% ===== 原文 pp.62–75（PDF p.77–90）=====
% 第 2 章：§2.1（末段，自 Definition 2.25 起）+ §2.2 B-模型（含 §2.2.1–§2.2.3）
% 章标题与 §2.1 标题由 c02a 写，本块不写；本块写 §2.2 标题。
% tag 区间：2.18–2.28（共 11 条，与 work/tags.txt 逐条对应）

% [原文 p.62]
\begin{definition}
设 $(\cM,\cE,\nabla)$ 是一个半无限 Hodge 结构变分（semi-infinite variation of Hodge structure）。若存在 $\cE$ 的一个截面 $s_0$，使得 $\cO_{\cM}$-模同态
\[
\cT_{\cM}\to\cE/\hbar\cE
\]
由
\[
X\longmapsto\hbar\nabla_Xs_0
\]
给出，并且是同构，则称该变分是极小 versal（miniversal）的。
\end{definition}

在极小 versal 的情形下，命题 2.24 连同一个附加的数据选择会在 $\cM$ 上诱导出一个 Frobenius 流形结构。更确切地说，假设给定如下数据：
\begin{itemize}[leftmargin=2em]
\item 一个半无限 Hodge 结构变分 $\cE,\nabla,(\cdot,\cdot)_{\cE}$，带有一个分次 $\mathrm{Gr}$。
\item 一个对置子空间（opposite subspace）$\cH_-\subseteq\cH$，它关于 $\overline{\Omega}$ 是迷向的，并且被 $\mathrm{Gr}$ 保持。
\item 一个元素 $\Omega_0\in\hbar\cH_-$，它代表 $[\Omega_0]\in\hbar\cH_-/\cH_-$，并且是 $\mathrm{Gr}_0$ 的一个特征向量。此外，相应的截面 $s'_0=\tau(\Omega_0\otimes1)$ 给出极小 versal 性。
\end{itemize}
于是，通过移动子空间构造（moving subspace construction）把 $\cE$ 视为 $\cH_{\cM}$ 的子丛时，我们照常把 $s'_0$ 认同为 $\cH_{\cM}$ 的一个截面，后者可以看作一个映射
\begin{equation}
\Psi:\cM\to\cH. \tag{2.18}
\end{equation}
这个映射可以如下描述。由于对 $q\in\cM$，$\cH$ 是直和 $\cH_-\oplus\cE_q$，故 $\Omega_0+\cH_-$ 与 $\cE_q$ 交于唯一一点，而该点事实上就是 $\Psi(q)$。这就是 Barannikov 的半无限 Hodge 结构变分的周期映射（period map）。我们于是可以定义映射
\begin{equation}
\psi:\cM\to\hbar\cH_-/\cH_- \tag{2.19}
\end{equation}
通过
\[
q\longmapsto[\hbar(\Psi(q)-\Omega_0)].
\]
极小 versal 性蕴含这是一个局部同构。事实上，在任意点 $q\in\cM$，我们需要验证 $\psi$ 的微分 $\psi_*:\cT_{\cM,q}\to\hbar\cH_-/\cH_-$ 是同构。现在 $\Psi$ 的微分是 $\Psi_*:\cT_{\cM,q}\to\cH$，它由 $X\mapsto\nabla_Xs'_0\in\cH_-$ 给出，这是因为，如命题 2.24 的 (1) 的证明中所论证的那样，$\cH_{\cM}$ 上的平凡联络限制为 $\cE$ 上的联络 $\nabla$。于是
\[
\psi_*(X)=[\hbar\nabla_Xs'_0]\in\hbar\cH_-/\cH_-\cong\cE_q/\hbar\cE_q.
\]
所以由极小 versal 性，$\psi_*$ 是同构。注意
\[
\psi_*(X)=A_X([\Omega_0]).
\]
这现在给出 $\cT_{\cM}$ 与平凡丛 $(\hbar\cH_-/\cH_-)_{\cM}$ 的一个等同，由
\[
X\longmapsto A_X([\Omega_0])
\]
给出。这允许我们把乘积 $(\cdot,\cdot)_{\hbar\cH_-/\cH_-}$ 搬到 $\cT_{\cM}$ 上，我们把它记作 $g$。此外，$(\hbar\cH_-/\cH_-)_{\cM}$ 上的平凡平坦联络给出 $\cT_{\cM}$ 上的平坦联络 $\nabla^{\cM}$。
我们通过条件
\[
A_{X\circ Y}[\Omega_0]=A_XA_Y[\Omega_0]
\]
来定义 $\cT_{\cM}$ 上的乘积，并定义张量
\[
\cA:S^3\cT_{\cM}\to\cO_{\cM}
\]
% [原文 p.63]
由
\[
\cA(X,Y,Z)=g(X\circ Y,Z)
\]
给出。它的对称性由 $\nabla$ 的平坦性推出：于是 $A_X$ 与 $A_Y$ 交换，并且用到式 (2.17)。设 $e$ 为 $\cM$ 上满足 $A_e[\Omega_0]=[\Omega_0]$ 的向量场。

\begin{theorem}
带数据 $(\nabla^{\cM},g,\cA)$ 以及向量场 $e$ 和 $E$ 的 $\cM$ 构成一个 Frobenius 流形，其中 $e$ 是平坦的单位元，而 $E$ 是一个 Euler 向量场，它满足式 (2.7)（常数 $d_0=1$）以及式 (2.6)（常数为 $2(\lambda+1)+D$，其中 $D$ 如定义 2.20 的 (7) 所述，而 $\lambda$ 是 $[\Omega_0]$ 关于 $\mathrm{Gr}_0$ 的特征值）。
\end{theorem}

\begin{proof}
显然 $(\nabla^{\cM},g,\cA)$ 在 $\cM$ 上定义了一个预 Frobenius 流形结构。同样显然 $e$ 是一个单位元；并且由于 $[\Omega_0]$ 是 $(\hbar\cH_-/\cH_-)_{\cM}$ 的常值、从而平坦的截面，$e$ 是平坦的。
由定理 2.11，只需证明由
\begin{align*}
\widehat{\nabla}^{\cM}_X(Y)&=\nabla^{\cM}_X(Y)+\hbar^{-1}X\circ Y,\\
\widehat{\nabla}_{\hbar\partial_\hbar}&=\hbar\partial_\hbar Y-\hbar^{-1}E\circ Y+\mathrm{Gr}_E(Y)
\end{align*}
定义的第一结构联络 $\widehat{\nabla}^{\cM}$ 是平坦的。
% [Tutor Note: p.63 第二个式子原文左端印作 $\widehat{\nabla}_{\hbar\partial_\hbar}$，漏掉了上标 $\cM$ 与括号内的 $Y$（第一个式子则完整写作 $\widehat{\nabla}^{\cM}_X(Y)$）；按上下文及 (2.20) 的写法，它应读作 $\widehat{\nabla}^{\cM}_{\hbar\partial_\hbar}(Y)$，即“算子 $\widehat{\nabla}^{\cM}_{\hbar\partial_\hbar}$ 作用在 $Y$ 上”。本译照原文逐符号抄录，不改动；PDF p.78 图像可核。]
我们来看看这些运算在 $\cE$ 上对应着什么。首先注意，在等同 $\psi_*:\cT_{\cM}\to(\hbar\cH_-/\cH_-)_{\cM}$ 之下，$\nabla^{\cM}$ 就是 $(\hbar\cH_-/\cH_-)_{\cM}$ 上的平坦联络，因此
\begin{align*}
(\tau\circ\psi_*)\bigl(\widehat{\nabla}^{\cM}_X(Y)\bigr)
&=\tau\Bigl(\psi_*\bigl(\nabla^{\cM}_X(Y)\bigr)+\hbar^{-1}\psi_*(X\circ Y)\Bigr)\\
&=\tau\Bigl(\dd\bigl(\psi_*(Y)\bigr)(X)+\hbar^{-1}A_{X\circ Y}([\Omega_0])\Bigr)\\
&=\tau\Bigl(\dd\bigl(A_Y([\Omega_0])\bigr)(X)+\hbar^{-1}A_XA_Y([\Omega_0])\Bigr)\\
&=\nabla_X\bigl((\tau\circ\psi_*)(Y)\bigr),
\end{align*}
这里最后一步用到命题 2.24 的 (1)。
其次，我们比较 $\mathrm{Gr}_E$ 与 $\mathrm{Gr}_0$。设 $Y$ 是平坦向量场，从而 $\mathrm{Gr}_E(Y)=[E,Y]$。则
\begin{align*}
\mathrm{Gr}_0(\psi_*(Y))&=\mathrm{Gr}_0(A_Y[\Omega_0])\\
&=\mathrm{Gr}(A_Y[\Omega_0])\quad\bmod\cH_-\\
&=\mathrm{Gr}(\hbar\nabla_Y\Omega_0)\quad\bmod\cH_-\\
&=\hbar\nabla_Y\Omega_0+\hbar\,\mathrm{Gr}(\nabla_Y\Omega_0)\quad\bmod\cH_-\\
&=\hbar\nabla_Y\Omega_0+\hbar\nabla_Y(\mathrm{Gr}\,\Omega_0)+\hbar\nabla_{[E,Y]}\Omega_0\quad\bmod\cH_-\\
&=A_Y[\Omega_0]+\lambda A_Y[\Omega_0]+A_{\mathrm{Gr}_E(Y)}[\Omega_0]\\
&=\psi_*\bigl((\lambda+1+\mathrm{Gr}_E)(Y)\bigr).
\end{align*}
这里第三个与第六个等号由命题 2.24 的 (1) 得到，第四个等号由定义 2.20 的 (5) 得到，第五个等号由定义 2.20 的 (6) 得到。

% [原文 p.64]
因此，对 $\cM$ 上的平坦向量场 $Y$ 以及 $f\in\cO_{\cM}\{\hbar\}$，可以计算
\begin{align*}
\tau\circ\psi_*\bigl(\widehat{\nabla}_{\hbar\partial_\hbar}(fY)\bigr)
&=\tau\circ\psi_*\Bigl((\hbar\partial_\hbar f)Y-\hbar^{-1}fE\circ Y+f\,\mathrm{Gr}_E(Y)\Bigr)\\
&=\tau\Bigl((\hbar\partial_\hbar f)\psi_*(Y)-\hbar^{-1}fA_E\psi_*(Y)+f(\mathrm{Gr}_0-\lambda-1)\psi_*(Y)\Bigr)\\
&=(\hbar\partial_\hbar f-f\nabla_E)(\tau\circ\psi_*)(Y)+\tau\bigl(\mathrm{Gr}_0(\psi_*(fY))\bigr)-(\lambda+1)\tau\circ\psi_*(fY)\\
&=(\hbar\partial_\hbar+E)(f)(\tau\circ\psi_*)(Y)+\tau\bigl(\mathrm{Gr}_0(\psi_*(fY))\bigr)-\nabla_E(\tau\circ\psi_*)(fY)\\
&\qquad-(\lambda+1)(\tau\circ\psi_*)(fY)\\
&=\mathrm{Gr}\bigl(\tau\circ\psi_*(fY)\bigr)-\nabla_E(\tau\circ\psi_*)(fY)-(\lambda+1)(\tau\circ\psi_*)(fY),
\end{align*}
最后一步由命题 2.24 的 (3) 得到。因此 $\widehat{\nabla}_{\hbar\partial_\hbar}$ 在 $\cE$ 上与 $\mathrm{Gr}-\nabla_E-(\lambda+1)$ 一致。
现在 $\nabla$ 是平坦联络，因此 $\widehat{\nabla}^{\cM}$ 在沿 $\cM$ 的切向方向上是平坦的。所以我们只需证明
\begin{equation}
[\widehat{\nabla}^{\cM}_{\hbar\partial_\hbar},\widehat{\nabla}^{\cM}_X]=0 \tag{2.20}
\end{equation}
对 $\cM$ 上的向量场 $X$ 成立。注意 $\nabla$ 的平坦性蕴含：对 $\cM$ 上的向量场 $X$ 与 $Y$，
\[
[\nabla_X,\nabla_Y]=\nabla_{[X,Y]}.
\]
把 (2.20) 搬到 $\cE$ 上，我们需要证明
\[
[\mathrm{Gr}-\nabla_E-(\lambda+1),\nabla_X]=0.
\]
由定义 2.20 的 (6)，计算左端给出
\[
[\mathrm{Gr},\nabla_X]-[\nabla_E,\nabla_X]=\nabla_{[E,X]}-\nabla_{[E,X]}=0.
\]
这说明 $\widehat{\nabla}^{\cM}$ 是平坦的，从而 $\cM$ 是 Frobenius 流形。
最后，我们核对常数 $D$ 的取值。对 $\cM$ 上的平坦向量场 $X$ 与 $Y$，有
\begin{align*}
&E(g(X,Y))-g([E,X],Y)-g(X,[E,Y])\\
&\quad=E\Bigl((\psi_*(X),\psi_*(Y))_{\hbar\cH_-/\cH_-}\Bigr)-(\psi_*(\mathrm{Gr}_E(X)),\psi_*(Y))_{\hbar\cH_-/\cH_-}\\
&\qquad-(\psi_*(X),\psi_*(\mathrm{Gr}_E(Y)))_{\hbar\cH_-/\cH_-}\\
&\quad=-\bigl((\mathrm{Gr}_0-(\lambda+1))\psi_*(X),\psi_*(Y)\bigr)_{\hbar\cH_-/\cH_-}\\
&\qquad-\bigl(\psi_*(X),(\mathrm{Gr}_0-(\lambda+1))\psi_*(Y)\bigr)_{\hbar\cH_-/\cH_-}\\
&\quad=(D+2(\lambda+1))(\psi_*(X),\psi_*(Y))_{\hbar\cH_-/\cH_-}\\
&\quad=(D+2(\lambda+1))g(X,Y).
\end{align*}
\end{proof}

注意如前所述，映射 $\psi$ 是 $\cM$ 与向量空间 $\hbar\cH_-/\cH_-$ 之间的局部同构，这给出 $\cM$ 上的一个线性结构。关于 $\nabla^{\cM}$ 的平坦向量场恰好是 $\hbar\cH_-/\cH_-$ 上的常值向量场。于是 $\hbar\cH_-/\cH_-$ 上的线性坐标诱导出 $\cM$ 上的坐标，其相伴的向量场是平坦的。我们称这些诱导坐标为平坦坐标；它们由上述数据的选择典范地给出。

% [Tutor Note: 定理 2.26（p.63）中的符号 $D$ 被原文重复使用：一句之内既说“(2.6) 的常数为 $2(\lambda+1)+D$”，又说“其中 $D$ 如定义 2.20 的 (7) 所述”，两处 $D$ 含义不同（前者是 (2.6) 的常数、后者是定义 2.20 的 $D$）。这是原文的记号重载，非排版错误；p.64 末行 $(D+2(\lambda+1))g(X,Y)$ 可作核对锚点，照译不改。]

% [原文 p.65]
\begin{example}
我们回到例 2.21 及其来自量子上同调的续篇例 2.22。空间 $\cH$ 认同为
\[
H^*(X,\C)\otimes_{\C}\C\{\hbar,\hbar^{-1}\},
\]
并取
\[
\cH_-:=H^*(X,\C)\otimes_{\C}\hbar^{-1}\cO(\PP^1\setminus\{0\})\subseteq\cH.
\]
回忆例 2.22：$\cE=H^*(X,\C)\otimes_{\C}\C\{\hbar\}$ 在 $\cH_{\widetilde{\cM}}$ 中的嵌入由映射 $\mathbb{J}$ 给出，它把 $\cE$ 的平坦截面送到 $\cH_{\widetilde{\cM}}$ 的常值截面。在 $\widetilde{\cM}$ 上坐标为 $y_0=\cdots=y_m=0$ 的点 $x$ 处，有
\[
\mathbb{J}(\alpha)=\alpha+O(\hbar^{-1}),
\]
从而 $\cE_x$ 的纤维在 $\mathbb{J}$ 下的像与 $\cH_-$ 只交于 $0$。因此
\[
\cH=\cH_-\oplus\cE_x.
\]
诱导出的同构
\[
\tau:(\hbar\cH_-/\cH_-)\otimes_{\C}\cO_{\widetilde{\cM}}\{\hbar\}\to\cE
\]
是一个映射
\[
\tau:H^*(X,\C)\otimes_{\C}\cO_{\widetilde{\cM}}\{\hbar\}\to\cE,
\]
它由
\[
\tau(s\otimes f)=f\cdot s'
\]
定义，其中 $s=s'+s''$，$s'$ 是 $\cE\subseteq\cH_{\widetilde{\cM}}$ 的截面，$s''$ 是 $\cH_{-,\widetilde{\cM}}$ 的截面。我们这样来做：
\[
s'=\mathbb{J}(s),\qquad s''=s-\mathbb{J}(s),
\]
注意由 $\mathbb{J}$ 的显式公式，$s-\mathbb{J}(s)$ 是 $\cH_{-,\widetilde{\cM}}$ 的截面，并且由于 $\mathbb{J}$ 给出 $\cE$ 在 $\cH_{\widetilde{\cM}}$ 中的嵌入，$\mathbb{J}(s)\in\cE\subseteq\cH_{\widetilde{\cM}}$。
为验证 $\cH_-$ 关于 $\overline{\Omega}$ 是迷向的，我们首先需要把 $\cE$ 上的配对延拓到 $\cH_-$ 的元素上。由于 $\mathbb{J}:\cE\otimes_{\C\{\hbar\}}\C\{\hbar,\hbar^{-1}\}\to\cH_{\widetilde{\cM}}$ 给出 $\cE\otimes_{\C\{\hbar\}}\C\{\hbar,\hbar^{-1}\}$ 与丛 $\cH_{\widetilde{\cM}}$ 的等同，我们把 $s_1,s_2\in\cH_{\widetilde{\cM}}$ 的 $(s_1,s_2)_{\cE}$ 解释为 $(\mathbb{J}^{-1}s_1,\mathbb{J}^{-1}s_2)_{\cE}$。但按定义 $\mathbb{J}^{-1}=S$，而由命题 2.23 的论证，$S^{-1}$ 是该配对下 $S$ 的伴随，所以
\[
(\mathbb{J}^{-1}s_1,\mathbb{J}^{-1}s_2)_{\cE}=(s_1,s_2)_{\cE}.
\]
于是显然，若 $s_1,s_2\in\cH_-$，则
\[
(s_1,s_2)_{\cE}=(\text{常数})\hbar^{-2}+\hbar^{-1}\text{ 的高阶项},
\]
从而 $\overline{\Omega}(s_1,s_2)=0$。因此 $\cH_-$ 关于 $\overline{\Omega}$ 是迷向的。
接着注意 $\cH_-$ 被 $\mathrm{Gr}$ 保持。同样，$\mathrm{Gr}$ 定义在 $\cE$ 上，所以给定 $s\in\cH_-$，我们应把 $\mathrm{Gr}(s)$ 解释为 $\mathbb{J}\circ\mathrm{Gr}\circ\mathbb{J}^{-1}(s)=\mathbb{J}\circ\mathrm{Gr}\circ S(s)$，由命题 2.18 以及 $\mathrm{Gr}$ 的定义，它位于 $\cH_-$ 中。此外，再次由命题 2.18，对 $s\in\hbar\cH_-$ 注意
\begin{equation}
\mathrm{Gr}(s)=\mathrm{Gr}_E(s)\quad\bmod\cH_-. \tag{2.21}
\end{equation}
特别地，命题 2.24 的 (1)–(3) 的假设得到满足。
我们现在将利用这一点，通过在 $\widetilde{\cM}$ 上选取 $[\Omega_0]=T_0\in\hbar\cH_-/\cH_-$（它由 $T_0\in\hbar\cH_-$ 代表）来构造一个 Frobenius 流形结构。我们需要验证：

% [原文 p.66]
\begin{itemize}[leftmargin=2em]
\item 相应的截面 $\mathbb{J}(T_0)\in\cE$ 给出极小 versal 性。注意 $\mathbb{J}(T_0)$ 正是该截面作为 $\cH_{\widetilde{\cM}}$ 的子丛 $\cE$ 的描述。作为抽象丛，$\cE\cong H^*(X,\C)\otimes_{\C}\cO_{\widetilde{\cM}}\{\hbar\}$，带有嵌入 $\mathbb{J}:\cE\to\cH_{\widetilde{\cM}}$。在这一等同下，$\cE$ 作为 $\cH_{\widetilde{\cM}}$ 的子丛的截面 $\mathbb{J}(T_0)$ 被认同为 $\cE$ 作为抽象丛的截面 $\mathbb{J}^{-1}(\mathbb{J}(T_0))=T_0$。于是
\[
\hbar\nabla_{\partial_{y_i}}(T_0)=\hbar\partial_{y_i}(T_0)+T_0*T_i=T_i.
\]
显然 $T_0,\dots,T_m$ 构成 $\cE/\hbar\cE$ 的一组基，于是我们得到极小 versal 性。
\item $[\Omega_0]$ 是 $\mathrm{Gr}_0$ 的一个特征向量。注意（第二个等号用到式 (2.21)）
\begin{align*}
\mathrm{Gr}_0[\Omega_0]&=\mathrm{Gr}(T_0)\quad\bmod\cH_-\\
&=\mathrm{Gr}_E(T_0)\quad\bmod\cH_-\\
&=-T_0\quad\bmod\cH_-\\
&=-[\Omega_0].
\end{align*}
所以 $[\Omega_0]$ 是特征值为 $-1$ 的特征向量。
\end{itemize}
于是所选数据在 $\widetilde{\cM}$ 上给出一个 Frobenius 流形结构。还要注意
\[
\psi:\widetilde{\cM}\to\hbar\cH_-/\cH_-
\]
由
\begin{equation}
\begin{aligned}
(y_0,\dots,y_m)\in\widetilde{\cM}\;&\longmapsto\;\hbar(\mathbb{J}(T_0)-T_0)\quad\bmod\cH_-\\
&=\sum_{i=0}^{m}y_iT_i
\end{aligned}
\tag{2.22}
\end{equation}
给出，这里用到命题 2.23 中 $\mathbb{J}$ 的公式。
我们来验证这个诱导出的 Frobenius 流形结构事实上就是量子上同调的 Frobenius 流形。确实，由 $A_X$ 的定义，$A_{\partial_{y_i}}T_j=T_i*T_j$。因此 $\psi_*$ 是恒等映射，乘法 $\circ$ 与 $*$ 给出的乘法一致，而联络 $\nabla^{\widetilde{\cM}}$ 是平坦联络。度量显然是正确的度量，$T_0$ 是单位元，Euler 向量场 $E$ 是标准的那个。
从上面的讨论我们看到，函数 $\mathbb{J}(T_0)$ 扮演了特殊的角色。这个函数有一个名字：
\end{example}

\begin{definition}
$X$ 的 Givental $J$-函数是
\[
J_X:=\mathbb{J}(T_0).
\]
\end{definition}

\begin{remark}
注意 $J_X$ 满足性质
\[
\hbar\partial_{y_i}J_X=\mathbb{J}(T_i).
\]
这可以通过直接计算来验证，也可以通过注意到，由于 $\mathbb{J}$ 给出 $\cE$ 在 $\cH_{\widetilde{\cM}}$ 中的嵌入，
\[
\mathbb{J}(T_i)=\mathbb{J}(T_0*T_i)=\mathbb{J}\bigl(\hbar\nabla_{\partial_{y_i}}T_0\bigr)=\hbar\partial_{y_i}\mathbb{J}(T_0)=\hbar\partial_{y_i}J_X.
\]
因此函数 $J_X$ 连同配对与分次算子，完全决定了半无限 Hodge 结构变分。
\end{remark}

% [Tutor Note: 例 2.27（p.65–66）中 $\mathbb{J}$ 是黑板体 J（原文用它表示 $H^*(X,\C)\otimes\cO_{\widetilde{\cM}}\{\hbar\}$ 的“平移/嵌入”映射，p.66 的 (2.22) 与定义 2.28 都用到它）；它与普通斜体 $J$（如 $J_X$）不是同一个符号。文字层 OCR 把 $\mathbb{J}$ 抽成“J”，须以 p.66 图像（PDF p.81）为准。]

% [原文 p.67]
\begin{example}
我们用命题 2.23 的公式写下 $X=\PP^2$ 的 Givental $J$-函数。对 $\beta\neq0$，基本类公理说明
\[
\langle T_0,T_2^m,\psi^\nu T_i\rangle_{0,\beta}=\langle T_2^m,\psi^{\nu-1}T_i\rangle_{0,\beta}.
\]
此外，由式 (2.2)，除非 $m+3\beta=2m+i+\nu-1$，即 $m=3\beta-i-\nu+1$，否则此式为零。另一方面，对 $\beta=0$，
\[
\langle T_0,T_2^m,\psi^\nu T_i\rangle_{0,0}=0
\]
除非 $m+1=2m+i+\nu$，即 $m=1-i-\nu$。但要使此式非零，$m$ 至少必须为 $1$，因为 $\overline{\cM}_{0,2}(X,0)$ 是空的，于是我们只在 $i=\nu=0$ 时得到贡献。特别地，$\langle T_0,T_2,T_0\rangle_{0,0}=\int_XT_0\cup T_2\cup T_0=1$。把这些一并代入命题 2.23 的公式并利用 $T^i=T_{2-i}$，我们得到
\[
\begin{aligned}
J_{\PP^2}={}&e^{(y_0T_0+y_1T_1)/\hbar}\cup\Biggl(T_0+\sum_{i=0}^{2}\biggl(y_2\hbar^{-1}\delta_{2,i}\\
&+\sum_{d\ge1}\sum_{\nu\ge0}\langle T_2^{3d+i-2-\nu},\psi^\nu T_{2-i}\rangle_{0,d}\hbar^{-(\nu+2)}e^{dy_1}\frac{y_2^{3d+i-2-\nu}}{(3d+i-2-\nu)!}\biggr)T_i\Biggr).
\end{aligned}
\]
\end{example}

% [Tutor Note: 例 2.30（p.67）中 $J_{\PP^2}$ 的公式原文没有编号，且其中“$\cup$”表示量子上同调中的 cup 积（不是普通乘法）；本译按图像逐符号抄录（PDF p.82）。另注意原文把 $\overline{\cM}_{0,2}(X,0)$ 印作 M 带上划线，指模空间 $\overline{\cM}_{0,2}(X,0)$，勿与正文中的 $\cM$ 混同。]

\subsection{B-模型}
我们现在考虑 $\PP^n$ 的镜像，并解释前一节所探讨的诸结构在 B-模型中的版本。
Calabi–Yau 流形的镜像仍是 Calabi–Yau 流形。另一方面，$\PP^n$ 是 Fano 流形，而 Fano 流形 $X$ 的镜像是所谓的 Landau–Ginzburg 模型。用数学语言说，这是一对 $(\check X,W)$，其中 $\check X$ 是一个簇，$W:\check X\to\C$ 是一个正则函数。这个正则函数被称为 Landau–Ginzburg 势，它自始至终起着关键作用。
环面 Fano 流形的镜像最先由 Givental 在 [33] 中描述，物理上的推导由 Hori 与 Vafa 在 [58] 中给出。一个关键点将是经由半无限 Hodge 结构变分这一中介，从数据 $(\check X,W)$ 构造 Frobenius 流形。此时镜像对称的预言是：这个 B-模型 Frobenius 流形与由 $X$ 的量子上同调产生的 Frobenius 流形重合。从 $(\check X,W)$ 构造 Frobenius 流形，在若干具体情形下已由 Sabbah [103] 与 Barannikov [4] 完成。这里我们将遵循 Barannikov 的途径。这应视为 Kyoji Saito 工作的推广：他在 [104] 中展示了如何把一个 Frobenius 流形结构联系到 $\C^n$ 上一个函数芽 $W$。
由于我们主要处理 $\PP^n$ 的情形、更具体地处理 $\PP^2$ 的情形，我们把这个情形下的镜像明确写出来，以便读者心中有一个例子。这一构造将在 §2.2.3 与第 5 章中更详细地探讨，但目前下面的做法就够了。
我们把 $\PP^n$ 的镜像 $\check X$ 描述为
\begin{equation}
\check X=(\C^\times)^n\subseteq\C^{n+1} \tag{2.23}
\end{equation}
由方程
\[
x_0\cdots x_n=1
\]
% [原文 p.68]
给出，其中 $x_0,\dots,x_n$ 是 $\C^{n+1}$ 上的坐标。于是势 $W$ 由
\[
W=x_0+\cdots+x_n
\]
给出。

\subsubsection{扭 de Rham 复形}
现在固定一个非奇异簇 $X$ 以及一个正则函数 $W:X\to\C$，即一个 Landau–Ginzburg 势。首先要考虑的问题是：与 $(X,W)$ 相关联的上同调群是什么？答案有几种形式。
考虑扭 de Rham 复形
\[
(\Omega^\bullet_X,\dd+\dd W\wedge).
\]
这是一个复形
\[
\Omega^0_X\to\Omega^1_X\to\Omega^2_X\to\cdots,
\]
其中微分为 $\omega\mapsto\dd\omega+\dd W\wedge\omega$。这里 $\Omega^p_X$ 是 $X$ 上 $p$-形式的层。对 $\Omega^p_X$ 这里有两种可能的解释：它可以是 Zariski 拓扑下的代数的 $p$-形式层，也可以是解析拓扑下的全纯 $p$-形式层。我们需要下面的定理，它最初由 Barannikov 与 Kontsevich 得到（未发表），随后由 Sabbah [102] 与 Ogus–Vologodsky [89] 给出。

\begin{theorem}
若 $W:X\to\C$ 是射影的，则
\[
\mathbb{H}^i_{\mathrm{Zar}}\bigl(X,(\Omega^\bullet_X,\dd+\dd W\wedge)\bigr)\cong\mathbb{H}^i_{\mathrm{Zar}}\bigl(X,(\Omega^\bullet_X,\dd W\wedge)\bigr)\cong\mathbb{H}^i_{\mathrm{An}}\bigl(X,(\Omega^\bullet_X,\dd+\dd W\wedge)\bigr).
\]
这里下标 Zar 与 An 表示我们处理的是 Zariski 拓扑下的代数版本或解析拓扑下的全纯版本。此外，$\mathbb{H}^i$ 表示超上同调。
\end{theorem}

我们所关心的 Landau–Ginzburg 势的例子，比如 $\PP^n$ 的镜像，事实上不是射影的，但事实证明这并不太重要：我们常常能找到部分紧化 $X\subseteq\overline{X}$，使得 $W$ 延拓为射影的 $W:\overline{X}\to\C$。尽管一般而言用 $X$ 与用 $\overline{X}$ 工作可能会有细微差别，但在我们所考虑的情形中不会有差别。我们在 §3.5 中详细讨论这一点。

\begin{example}
假设处于上述定理的情形，并且 $W$ 只有孤立临界点。那么用超上同调谱序列
\[
E_2^{p,q}=H^p\bigl(X,\mathcal{H}^q(\Omega^\bullet,\dd W\wedge)\bigr)\Rightarrow\mathbb{H}^n\bigl(X,(\Omega^\bullet_X,\dd W\wedge)\bigr)
\]
不难计算 $\mathbb{H}^i(X,(\Omega^\bullet_X,\dd W\wedge))$。这里 $H^p$ 是层上同调，而 $\mathcal{H}^q$ 表示复形 $(\Omega^\bullet_X,\dd W\wedge)$ 的上同调。
首先，我们来计算 $\mathcal{H}^q(\Omega^\bullet_X,\dd W\wedge)$。多重线性代数的一个基本事实是：若 $V$ 是有限维向量空间且 $v\in V$ 非零，则复形
\[
\textstyle\bigwedge^0V\xrightarrow{\wedge v}\bigwedge^1V\xrightarrow{\wedge v}\bigwedge^2V\xrightarrow{\wedge v}\cdots
\]
正合。由于 $\dd W=0$ 恰好成立在 $W$ 的临界轨迹上（我们把它记作 $\mathrm{Crit}(W)$），层 $\mathcal{H}^q(\Omega^\bullet_X,\dd W\wedge)$ 支撑在 $\mathrm{Crit}(W)$ 上。此外，若 $q<\dim X=:n$，我们将归纳地证明 $\mathcal{H}^q(\Omega^\bullet_X,\dd W\wedge)=0$。当然 $\mathcal{H}^0(\Omega^\bullet_X,\dd W\wedge)=0$，因为它一方面支撑在 $\mathrm{Crit}(W)$ 上，另一方面包含在局部自由层 $\Omega^0_X$ 中。
% [原文 p.69]
若对 $q'<q$ 有 $\mathcal{H}^{q'}(\Omega^\bullet_X,\dd W\wedge)=0$，则我们得到正合列
\[
0\to\Omega^0_X\xrightarrow{\dd W\wedge}\cdots\xrightarrow{\dd W\wedge}\Omega^{q-1}_X\to\ker(\Omega^q_X\to\Omega^{q+1}_X)\to\mathcal{H}^q(\Omega^\bullet_X,\dd W\wedge)\to0.
\]
现在令 $\mathcal{F}=\operatorname{coker}\bigl(\Omega^{q-2}_X\xrightarrow{\dd W\wedge}\Omega^{q-1}_X\bigr)$；层 $\mathcal{F}$ 的投射维数至多为 $q-1$，因此深度至少为 $n-(q-1)>1$（见 [57], Proposition III 6.12A）。另一方面，$\ker(\Omega^q_X\to\Omega^{q+1}_X)$ 的深度至少为 $1$，因为它是无挠的。但这样一来，只要 $\mathcal{H}^q(\Omega^\bullet_X,\dd W\wedge)$ 非零，$\mathcal{H}^q(\Omega^\bullet_X,\dd W\wedge)$ 的深度就至少为 $1$（见 [57], III Exercise 3.4），这与它支撑在零维集合 $\mathrm{Crit}(W)$ 上相矛盾。因此 $\mathcal{H}^q(\Omega^\bullet_X,\dd W\wedge)=0$。
接下来考察 $\mathcal{H}^n(\Omega^\bullet_X,\dd W\wedge)$。若我们在 $X$ 的临界点附近有局部坐标 $x_1,\dots,x_n$，则有 $\Omega^{n-1}_X\to\Omega^n_X$ 由
\[
\dd x_1\wedge\cdots\wedge\widehat{\dd x_i}\wedge\cdots\wedge\dd x_n\longmapsto\pm\frac{\partial W}{\partial x_i}\dd x_1\wedge\cdots\wedge\dd x_n
\]
给出。因此这个映射的像局部地恰好是 $W$ 的雅可比理想。
由于 $\mathcal{H}^q(\Omega^\bullet_X,\dd W\wedge)$ 现在只在 $q=n$ 时非零，而对 $q=n$ 该层的支撑是零维的，谱序列在 $E_2$ 项退化，于是我们得到
\[
\mathbb{H}^p\bigl(X,(\Omega^\bullet_X,\dd W\wedge)\bigr)=
\begin{cases}
0,&p\neq n,\\[2pt]
\Gamma\Bigl(X,\operatorname{coker}\bigl(\Omega^{n-1}_X\xrightarrow{\dd W\wedge}\Omega^n_X\bigr)\Bigr),&p=n.
\end{cases}
\]
\hfill$\square$
\end{example}

\begin{example}
设 $(X,W)$ 是 (2.23) 中给出的 $\PP^n$ 的镜像，并作一点小改动。我们实际上考虑 Landau–Ginzburg 势的一种可能的变动，即令 $x_0\cdots x_n=\kappa$，其中 $\kappa\in\C^\times$，而不是 $x_0\cdots x_n=1$。这给出以 $\kappa$ 为参数的单参数 Landau–Ginzburg 模型族。当然，如前所述，为应用上述理论我们需要把 $X$ 部分紧化，但事实上这不会增加任何临界点，所以计算是一样的。关于这一部分紧化与计算的更多细节见 §3.5。因此，不做部分紧化，我们通过 $n$-形式
\[
\Omega=\frac{\dd x_1\wedge\cdots\wedge\dd x_n}{x_1\cdots x_n}
\]
把 $\Omega^n_X$ 平凡化，然后
\[
\mathbb{H}^n\bigl(X,(\Omega^\bullet_X,\dd W\wedge)\bigr)=\frac{\C[x_1^{\pm},\dots,x_n^{\pm}]}{(\partial W/\partial x_1,\dots,\partial W/\partial x_n)}\cdot\Omega.
\]
这里我们消去了变量 $x_0=\kappa/(x_1\cdots x_n)$。环 $\C[x_1^{\pm},\dots,x_n^{\pm}]/\mathrm{Jac}(W)$（其中 $\mathrm{Jac}(W)=(\partial W/\partial x_1,\dots,\partial W/\partial x_n)$ 是雅可比理想）称为 $W$ 的 Milnor 环。
它是一个有限维 $\C$-代数；我们来计算它。用雅可比理想的生成集 $x_i\partial W/\partial x_i$ 更为方便，注意
\[
x_i\frac{\partial W}{\partial x_i}=x_i-\frac{\kappa}{x_1\cdots x_n}.
\]
于是 $W$ 的临界点出现在
\[
x_1=\cdots=x_n=\frac{\kappa}{x_1\cdots x_n},
% [原文 p.70]
\]
即 $(x_1,\dots,x_n)=(\mu,\dots,\mu)$，其中 $\mu=\mu^{-n}\kappa$，亦即 $\mu$ 是 $\kappa$ 的一个 $(n+1)$ 次单位根。此外，Milnor 环没有幂零元。所以它作为 $\C$-向量空间的维数是 $n+1$，与 $H^*(\PP^n,\C)$ 的维数相同。\hfill$\square$
\end{example}

% [Tutor Note: 例 2.32、2.33、2.34（p.69–72）原文都在例内附有计算并以 $\square$ 封口，这不是行文错误：本书把带证明的计算放进 Example 环境。另 p.70 的 $\mu=\mu^{-n}\kappa$ 是原文写法，与 $\mu^{n+1}=\kappa$（$\mu$ 为 $\kappa$ 的 $(n+1)$ 次单位根）等价，照译。]

接着回到扭 de Rham 复形，我们注意到可以修改这一扭化，得到一族以 $\hbar\in\C^\times$ 为参数的扭 de Rham 复形。我们取复形 $(\Omega^\bullet_X,\dd+\hbar^{-1}\dd W\wedge)$。当然，$\mathbb{H}^p(X,(\Omega^\bullet_X,\dd+\hbar^{-1}\dd W\wedge))$ 的维数不依赖于 $\hbar$，因为复形 $(\Omega^\bullet_X,\hbar^{-1}\dd W\wedge)$ 同构于 $(\Omega^\bullet_X,\dd W\wedge)$。

\subsubsection{同调}
给定一个 Landau–Ginzburg 模型 $(X,W)$，我们得到一个自然的同调理论，它与扭 de Rham 复形的上同调对偶。为简化讨论，我们如例 2.32 那样假设 $\mathrm{Crit}(W)$ 是零维的，从而只有 $\mathbb{H}^n(X,(\Omega^\bullet_X,\dd+\hbar^{-1}\dd W\wedge))$ 非零，其中 $n=\dim X$。此外，若 $X$ 是仿射的，$\mathbb{H}^n(X,(\Omega^\bullet_X,\dd+\hbar^{-1}\dd W\wedge))$ 的任一元素都由一个 $n$-形式代表。确实，我们可以通过第二个超上同调谱序列看到这一点：
\[
E_1^{p,q}=H^p\bigl(H^q(X,\Omega^\bullet_X),\dd+\hbar^{-1}\dd W\wedge\bigr)\Rightarrow\mathbb{H}^n\bigl(X,(\Omega^\bullet_X,\dd+\hbar^{-1}\dd W\wedge)\bigr),
\]
其中我们对复形 $\Omega^\bullet_X$ 的每个条目取 $q$ 次层上同调，然后对所得复形取 $p$ 次上同调。由于 $X$ 是仿射的，对 $q>0$ 有 $H^q(X,\Omega^r_X)=0$，从而谱序列退化，这表明
\[
\mathbb{H}^n\bigl(X,(\Omega^\bullet_X,\dd+\hbar^{-1}\dd W\wedge)\bigr)\cong H^n\bigl(\Gamma(X,\Omega^\bullet_X),\dd+\hbar^{-1}\dd W\wedge\bigr).
\]
在这种情形下，对固定的 $\hbar\in\C^\times$，我们考虑一个同调群，把它记作
\[
H_n(X,\operatorname{Re}W/\hbar\ll0;\C).
\]
我们指的是允许无界闭链的同调理论，但只在那些使 $\operatorname{Re}W/\hbar\to-\infty$ 的方向上无界。我们在这里不做得更精确；精确化是可能的，但对我们来说这只是一个技术点。重要的是：
\[
H_n(X,\operatorname{Re}W/\hbar\ll0;\C)
\]
自然地与
\[
\mathbb{H}^n\bigl(X,(\Omega^\bullet_X,\dd+\hbar^{-1}\dd W\wedge)\bigr)
\]
对偶。确实，给定一个闭链 $\Xi\in H_n(X,\operatorname{Re}W/\hbar\ll0;\C)$，以及代数地代表 $\mathbb{H}^n(X,(\Omega^\bullet_X,\dd+\hbar^{-1}\dd W\wedge))$ 中元素的 $\omega\in H^0(X,\Omega^n_X)$，我们可以计算所谓的振荡积分（oscillatory integral）
\[
\int_\Xi e^{W/\hbar}\omega.
\]
在 $H_n(X,\operatorname{Re}W/\hbar\ll0;\C)$ 的恰当定义下，这个积分收敛，正是因为沿 $\Xi$ 的无界方向有 $\operatorname{Re}W/\hbar\to-\infty$。此外注意，若 $\eta$ 是 $X$ 上的 $(n-1)$-形式，则
\[
\dd(e^{W/\hbar}\eta)=e^{W/\hbar}(\dd\eta+\hbar^{-1}\dd W\wedge\eta).
\]
由 Stokes 定理，
\[
\int_\Xi e^{W/\hbar}\omega=\int_\Xi e^{W/\hbar}\omega+\dd(e^{W/\hbar}\eta),
% [原文 p.71]
\]
% [Tutor Note: p.70 末的 Stokes 公式右端省略了括号：原文写作 $\int_\Xi e^{W/\hbar}\omega+\dd(e^{W/\hbar}\eta)$，应理解为 $\int_\Xi\bigl(e^{W/\hbar}\omega+\dd(e^{W/\hbar}\eta)\bigr)$（Stokes 定理把 $\int_\Xi\dd(e^{W/\hbar}\eta)$ 化为边界积分，$\Xi$ 为闭链时边界项为零）。照译不改；核对锚点是 p.71 开头“于是 $\omega$ 与 $\omega+(\dd+\hbar^{-1}\dd W\wedge)\eta$ 给出相同的积分”。]
于是 $\omega$ 与 $\omega+(\dd+\hbar^{-1}\dd W\wedge)\eta$ 给出相同的积分。它们也代表 $\mathbb{H}^n(X,(\Omega^\bullet_X,\dd+\hbar^{-1}\dd W\wedge))$ 中的同一个元素，于是我们得到后来被证明是自然的完美配对
\begin{equation}
H_n(X,\operatorname{Re}W/\hbar\ll0;\C)\times\mathbb{H}^n\bigl(X,(\Omega^\bullet_X,\dd+\hbar^{-1}\dd W\wedge)\bigr)\to\C. \tag{2.24}
\end{equation}
此外还有一个自然的相交配对
\begin{equation}
H_n(X,\operatorname{Re}W/\hbar\ll0;\C)\times H_n\bigl(X,\operatorname{Re}W/(-\hbar)\ll0;\C\bigr)\to\C. \tag{2.25}
\end{equation}
粗略地说，左边第一个与第二个同调群中的两个闭链只能在 $X$ 的某个有界区域内相交，且相交数是良定义的，因为这些闭链无法形变到使交点跑到无穷远。

\begin{example}
我们考虑 $\PP^2$ 的镜像，取 $(\check X,W)=((\C^\times)^2,W)$。这里 $(\C^\times)^2=V(x_0x_1x_2-1)\subseteq\C^3$，且 $W=x_0+x_1+x_2$。如我们在例 2.33 中所见，$\mathbb{H}^2(\check X,(\Omega^\bullet_{\check X},\dd+\hbar\dd W\wedge))$ 是三维的。我们将描述 $H_2(\check X,\operatorname{Re}W/\hbar\ll0;\C)$ 的一组基 $\Xi_0,\Xi_1,\Xi_2$。
为描述这组基，我们通过映射
\[
\operatorname{Log}:(\C^\times)^2\to\R^2,\qquad \operatorname{Log}(x_1,x_2)=(\log|x_1|,\log|x_2|)
\]
把 $(\C^\times)^2$ 认同为 $\R^2$ 上的平凡 $T^2$-丛。我们取 $\Xi_0=\operatorname{Log}^{-1}(0,0)$。这是一个紧闭链，所以对任意 $\hbar$，它当然代表 $H_2(\check X,\operatorname{Re}W/\hbar\ll0;\C)$ 中的一个类。然而，要使它真正成为一个闭链，我们必须选择定向。利用映射 $\operatorname{Log}$，我们可以把 $(\C^\times)^2$ 认同为 $\R^2\times T^2$，取 $\R^2$ 上的坐标 $y_1,y_2$ 与 $T^2$ 上的坐标 $\theta_1,\theta_2$，使得 $x_j=\exp(y_j+\ii\theta_j)$。于是我们可以用 $\dd\theta_1\wedge\dd\theta_2$ 给 $\Xi_0$ 定向。
接着我们定义 $\Xi_1$。设 $\rho_0,\rho_1,\rho_2\subseteq\R^2$ 为例 1.14 中扇的一维锥。令
\begin{align*}
S_0&=\bigl\{(x_1,x_2)\in\operatorname{Log}^{-1}(\rho_0)\bigm|\arg(x_1^{-1}x_2^{-1})=\arg(\hbar)+\pi\bigr\},\\
S_1&=\bigl\{(x_1,x_2)\in\operatorname{Log}^{-1}(\rho_1)\bigm|\arg x_1=\arg(\hbar)+\pi\bigr\},\\
S_2&=\bigl\{(x_1,x_2)\in\operatorname{Log}^{-1}(\rho_2)\bigm|\arg x_2=\arg(\hbar)+\pi\bigr\}.
\end{align*}
每个 $S_i$ 是一个柱面，其边界落在 $\operatorname{Log}^{-1}(0,0)$ 上，我们在图 2 中把它画成一个对边相互等同的正方形。我们取 $S_3$ 为包含在 $\operatorname{Log}^{-1}(0,0)$ 中、以 $\partial S_0\cup\partial S_1\cup\partial S_2$ 为边界的一个曲面，并取 $\Xi_1$ 为分片光滑闭链
\[
\Xi_1=S_0\cup S_1\cup S_2\cup S_3.
\]
比如说，我们用 $-\dd y_1\wedge\dd\theta_2$ 给 $S_1\subseteq\Xi_1$ 定向。可以验证这能延拓为整个 $\Xi_1$ 上的一个定向。
例如，图 2 画出了 $\hbar=1$ 时这样的一个曲面 $S_3$。为验证 $\Xi_1$ 位于 $H_2(\check X,\operatorname{Re}W/\hbar\ll0;\C)$ 中，注意在 $S_0$ 上，当沿 $S_0$ 的非紧方向前进时，$|x_1|$ 与 $|x_2|$ 趋于零，而 $|x_0|=1/|x_1x_2|\to\infty$。由于
\[
\arg(x_1^{-1}x_2^{-1})=\arg\hbar+\pi,
\]
事实上 $x_1^{-1}x_2^{-1}/\hbar$ 是实数并且趋于 $-\infty$。于是这一项在 $W/\hbar$ 中占主导，所以 $\operatorname{Re}W/\hbar\to-\infty$。同样的论证对 $S_1$ 与 $S_2$ 也成立。因此 $\Xi_1$ 位于 $H_2(\check X,\operatorname{Re}W/\hbar\ll0;\C)$ 中。

% [Tutor Note: 例 2.34（p.71）原文写作 $\mathbb{H}^2(\check X,(\Omega^\bullet_{\check X},\dd+\hbar\dd W\wedge))$，用的是 $\hbar$（一次幂）；而本节其余各处（如 p.70、p.71 的 (2.24)）均为 $\hbar^{-1}$。此处疑为原文笔误（应为 $\dd+\hbar^{-1}\dd W\wedge$），照译不改，见 PDF p.86 图像。]

% [原文 p.72]
\begin{figure}[htbp]
\centering
\includegraphics[width=0.42\textwidth]{f2_2.png}
\figtitle{2}{包含在二维环面中的曲面 $S_3$。正方形的对边相互等同。}
\end{figure}
% [Tutor Note: 图 2 按 STYLE_GUIDE §5 从 PDF p.87（原书 p.72）裁出，存为 work/figures/f2_2.png，未裁入原图注；图题按原文译为中文。原图为二维环面基本域（正方形对边等同）中 $S_3$ 的两块阴影区域。]

注意我们对 $S_3$ 有一个任意的选择，而它总可以通过加上 $\operatorname{Log}^{-1}(0,0)$ 的若干倍来修改。事实上，如果我们作一个选择，比如对 $\hbar=1$ 作选择，并让 $\hbar$ 沿 $\hbar=e^{2\pi\ii t}$、$0\le t\le1$ 连续变化，我们可以连续地选取 $S_3$。当绕行一圈回到起点时，我们会发现 $\Xi_1$ 被 $\Xi_1-3\Xi_0$ 取代（如果按照下面指出的方式给 $\Xi_0$ 与 $\Xi_1$ 定向）。这说明 $S_3$ 没有典范的选择。
为避免复杂性，我们只对 $\hbar=-1$ 描述 $\Xi_2$。取
\[
\Xi_2=\bigl\{(x_1,x_2)\in\check X\bigm|x_1,x_2\text{ 为实数},\ x_1,x_2>0\bigr\}.
\]
那么在非紧方向上同样有 $\operatorname{Re}W/\hbar\to-\infty$。我们用 $\dd y_1\wedge\dd y_2$ 给 $\Xi_2$ 定向。
我们把构造一个如 $\Xi_1$ 的例子中那样依赖于 $\hbar$ 的 $\Xi_2$ 留给读者。但要注意，我们所描述的这个 $\Xi_2$ 实际上对任何满足 $\operatorname{Re}\hbar<0$ 的 $\hbar$ 都适用。
我们也可以计算 (2.25) 中配对下的相交数。我们只记下最简单的一个：闭链 $\Xi_0$ 不依赖于 $\hbar$，并且是 $\operatorname{Log}$ 的一条纤维。当然任意两条纤维是同调的。所以 $\Xi_0\cdot\Xi_0=0$，且对任意 $\hbar$ 有 $\Xi_0\cdot\Xi_1(\hbar)=0$，因为 $\Xi_0$ 可以移到与 $\Xi_1(\hbar)$ 不相交。（这里我们标出 $\Xi_1$ 对 $\hbar$ 的依赖。）另一方面，$\Xi_0$ 与 $\Xi_2$ 横截地交于一点，而在所给定向选择下，$\Xi_0\cdot\Xi_2=-1$，这里我们把 $\Xi_0\in H_2(\check X,\operatorname{Re}W\ll0;\C)$ 与 $\Xi_2\in H_2(\check X,\operatorname{Re}(-W)\ll0;\C)$ 这样看待。\hfill$\square$
\end{example}

上面例子中给出的闭链描述并不是通常的描述；我们之所以这样给出，是因为这些闭链与 $H^*(\PP^2,\C)$ 的 $T_0,T_1,T_2$ 的镜像闭链联系更紧密（但并不相同）。$H_n(X,\operatorname{Re}W/\hbar\ll0;\C)$ 的生成元的更通常的描述是用 Lefschetz 针状体给出的。
假设 $W$ 的临界轨迹是约化的：临界轨迹通过 Milnor 环的谱被赋予概形结构。那么每个临界点都是非退化的，意思是若 $X$ 是 $n$ 维的，则在每个临界点附近存在局部全纯坐标 $z_1,\dots,z_n$，使得 $W=z_1^2+\cdots+z_n^2$。例如，对于 $\PP^n$ 的镜像 $(\check X,W)$ 就是这种情形。
在 $X$ 上放一个度量，并把函数 $\operatorname{Re}W/\hbar$ 视为 $X$ 上的 Morse 函数。在每个临界点附近局部地，存在实坐标 $y_1,\dots,y_{2n}$，使得
% [原文 p.73]
$\operatorname{Re}W/\hbar=y_1^2+\cdots+y_n^2-y_{n+1}^2-\cdots-y_{2n}^2$。所以每个临界点都是指标为 $n$ 的临界点。临界点 $p$ 的稳定流形 $\Delta_p^+$ 是 $\operatorname{Re}W/\hbar$ 的、当时间 $t\to+\infty$ 时趋于 $p$ 的梯度流线的并；临界点 $p$ 的不稳定流形 $\Delta_p^-$ 是 $\operatorname{Re}W/\hbar$ 的、当时间 $t\to-\infty$ 时趋于 $p$ 的梯度流线的并。当然，$\Delta_p^+$ 与 $\Delta_p^-$ 依赖于 $\hbar$，而在本讨论中我们认为它是固定的。
$\Delta_p^+$ 与 $\Delta_p^-$ 都是 $n$ 维子流形，并且若 $W$ 是真映射，则显然
\[
\Delta_p^+\in H_n(X,\operatorname{Re}W/\hbar\ll0;\C)
\]
且
\[
\Delta_p^-\in H_n(X,\operatorname{Re}(-W/\hbar)\ll0;\C).
\]
若 $W$ 不是真映射，我们可以取紧合化（properification）$X\subseteq\overline{X}$、$W:\overline{X}\to\C$ 真映射，并选取度量使得当 $p\in\mathrm{Crit}(W)$ 含于 $X$ 时 $\Delta_p^{\pm}\subseteq X$。所以我们仍得到
\[
\Delta_p^{\pm}\in H_n(X,\operatorname{Re}(\pm W/\hbar)\ll0;\C).
\]
这些闭链被称为 Lefschetz 针状体（Lefschetz thimble）。
对一般的度量选择，除非 $p=p'$，有 $\Delta_p^{\pm}\cap\Delta_{p'}^{\pm}=\varnothing$，且 $\Delta_p^+\cap\Delta_p^-=\{p\}$。于是取恰当的定向，我们可以假设在相交配对 (2.25) 下
\begin{equation}
\Delta_p^+\cdot\Delta_{p'}^-=\delta_{pp'}. \tag{2.26}
\end{equation}
我们将需要下面这个关于这些闭链上振荡积分的标准近似，它描述当 $\hbar\to0$ 时这些积分的渐近行为。这一近似被称为平稳相位近似（stationary phase approximation）：

\begin{proposition}
设 $X,W$ 如上，且 $W$ 只有孤立的非退化临界点。设 $f(x,\hbar)$ 是
\[
\check X\times\{|\hbar|<\epsilon\}\subseteq\check X\times\C
\]
上的全纯函数，并且对任意满足 $|\hbar|<\epsilon$ 的 $\hbar$，限制到 $\check X\times\{\hbar\}$ 上时是代数的。则对 $W$ 的临界点 $p$，
\[
\int_{\Delta_p^+}f(x,\hbar)e^{W/\hbar}\omega=\frac{(-2\pi\hbar)^{n/2}}{\sqrt{\mathrm{Hess}(W)(p)}}e^{W(p)/\hbar}\bigl(f(p,0)+O(\hbar)\bigr),
\]
其中 $\omega$ 是 $X$ 上处处非零的代数 $n$-形式，$n=\dim X$，$f$ 是正则函数，而 $p$ 处的 Hessian $\mathrm{Hess}(W)$ 是 $\det(\partial^2W/\partial x_i\partial x_j)$ 在 $p$ 处的值，这里 $x_1,\dots,x_n$ 是满足 $\omega=\dd x_1\wedge\cdots\wedge\dd x_n$ 的局部坐标。
\end{proposition}

\begin{proof}[（略证）]
首先注意，我们可以写
\begin{align*}
\int_{\Delta_p^+}f(x,\hbar)e^{W/\hbar}\omega&=\int_{\Delta_p^+}f(x,\hbar)e^{(W-W(p))/\hbar}e^{W(p)/\hbar}\omega\\
&=e^{W(p)/\hbar}\int_{\Delta_p^+}f(x,\hbar)e^{(W-W(p))/\hbar}\omega,
\end{align*}
于是我们可以假设 $W(p)=0$，代价是乘上因子 $e^{W(p)/\hbar}$。现在的要点是，当 $\hbar\to0$ 时的渐近行为来自积分
\begin{equation}
\int_{\Delta_p^+\cap B_\epsilon}f(x,\hbar)e^{W/\hbar}\omega, \tag{2.27}
\end{equation}
% [原文 p.74]
其中 $B_\epsilon$ 是以 $p$ 为球心、半径为 $\epsilon$ 的球。这是因为
\begin{equation}
\int_{\Delta_p^+\setminus B_\epsilon}f(x,\hbar)e^{W/\hbar}\omega\to0 \tag{2.28}
\end{equation}
当 $\hbar\to0$，因为 $\operatorname{Re}W/\hbar$ 在 $\Delta_p^+\setminus B_\epsilon$ 上有界地远离零且为负。于是 $e^{W/\hbar}$ 当 $\hbar\to0$ 时快速趋于 $0$，故 (2.28) 成立。积分 (2.27) 然后可以通过把 $f(x,\hbar)$ 换成它在 $x=p,\hbar=0$ 处的值、并把 $W$ 换成其 Taylor 展开中的二阶（首）项来近似：
\[
W=\sum_{i,j}\frac12\frac{\partial^2W}{\partial x_i\partial x_j}x_ix_j+\text{高阶项}.
\]
令 $A=((\partial^2W/\partial x_i\partial x_j)(p))_{i,j}$，于是 $\mathrm{Hess}(W)(p)=\det A$，我们可以作坐标变换 $y=\sqrt{A}x$（取 $A$ 的一个平方根），使 $W$ 的二阶部分现在是 $(y_1^2+\cdots+y_n^2)/2$，并且在 $p$ 附近局部地
\[
\omega=\dd y_1\wedge\cdots\wedge\dd y_n/\sqrt{\mathrm{Hess}(W)(p)}.
\]
现在取 $\Delta_p^+$ 对 $\hbar=1$ 定义；然后我们可以让 $\hbar$ 在 $\hbar$-平面中一个包含 $1$ 的角扇区内趋于 $0$。若写 $y_j=u_j+\ii v_j$，则
\[
\operatorname{Re}W=\frac12(u_1^2+\cdots+u_n^2-v_1^2-\cdots-v_n^2),
\]
所以 $\operatorname{Re}W$ 的稳定流形由 $u_1=\cdots=u_n=0$ 给出。于是，在这一角扇区内，我们可以用标准 Gauss 积分来近似 (2.27)：
\begin{align*}
&\frac{f(p,0)}{\sqrt{\mathrm{Hess}(W)(p)}}\int_{\R^n}e^{-(v_1^2+\cdots+v_n^2)/2\hbar}\dd(\ii v_1)\wedge\cdots\wedge\dd(\ii v_n)\\
&\quad=\frac{f(p,0)}{\sqrt{\mathrm{Hess}(W)(p)}}\sqrt{(2\pi\hbar)^n}\,\ii^n\\
&\quad=f(p,0)\frac{(-2\pi\hbar)^{n/2}}{\sqrt{\mathrm{Hess}(W)(p)}}.
\end{align*}
这给出所需的渐近行为。
\end{proof}

\subsubsection{B-模型的半无限 Hodge 结构变分}
我们首先需要讨论 B-模型的半无限 Hodge 结构变分所生活的 B-模型模空间。在 $\PP^n$ 的镜像的情形，我们从 $(\check X,W_0)$ 出发，其中 $W_0=x_0+\cdots+x_n$，$\check X=V(x_0\cdots x_n-1)\subseteq\C^{n+1}$。B-模型模空间是 $W_0$ 的万有开折（universal unfolding）。在复流形的语境中，万有开折是一个复流形芽 $0\in\cM$，连同全纯函数
\[
W:\cM\times\check X\to\C,
\]
使得：
\begin{itemize}[leftmargin=2em]
\item $W|_{\{0\}\times\check X}=W_0$；
\item 对任意复空间芽 $0\in Y$，以及满足 $W'|_{\{0\}\times\check X}=W_0$ 的全纯函数 $W':Y\times\check X\to\C$，
% [原文 p.75]
存在交换图表
\[
\begin{array}{ccc}
Y\times\check X&\xrightarrow{\ \eta\ }&\cM\times\check X\\[4pt]
\big\downarrow&&\big\downarrow\\[4pt]
Y&\xrightarrow{\ \varphi\ }&\cM
\end{array}
\]
其中 $\varphi:(0\in Y)\to(0\in\cM)$ 是芽的映射，且其在 Zariski 切空间上诱导的映射唯一，竖直箭头是投影，$\eta|_{\{0\}\times\check X}:\{0\}\times\check X\to\{0\}\times\check X$ 是恒等映射，并且 $W'=W\circ\eta$。
\end{itemize}
事实上，一个标准事实（例如见 [109]）说：假设 $W_0$ 有孤立临界点，则 $W_0$ 的万有开折是 Milnor 环
\[
\C[x_1^{\pm},\dots,x_n^{\pm}]/(\partial W_0/\partial x_1,\dots,\partial W_0/\partial x_n)
\]
中 $0$ 的一个芽，把它看成一个仿射空间。若 $f_0,\dots,f_n$ 是 $\C[x_1^{\pm},\dots,x_n^{\pm}]$ 中代表该向量空间一组基的元素，给出原点的芽 $\cM$ 上的坐标 $t_0,\dots,t_n$，则 $W_0$ 在 $\cM\times\check X$ 上的万有开折 $W$ 的方程为
\[
W:=W_0+\sum_{i=0}^{n}t_if_i.
\]
例如，我们可以取基 $f_i=W_0^i$，$0\le i\le n$。我们把验证这确实是 Milnor 环的一组基留给读者。
理想情况下，我们希望 $\cM$ 就是 B-模型的半无限 Hodge 结构变分所生活的空间。然而，当 $n\ge2$ 时会出现一个问题：当 $i\ge2$ 的 $t_i\neq0$ 时，$W_0+\sum t_iW_0^i$ 比 $W_0$ 有更多的临界点。例如，考虑 $n=2$ 的情形以及 $W_0$ 的小扰动
\[
W=W_0+t_2W_0^2=W_0(1+t_2W_0),
\]
其中 $t_2$ 很小。则 $W$ 分解为
\[
\check X\xrightarrow{W_0}\C\xrightarrow{\ f\ }\C,
\]
其中 $f(x)=x(1+t_2x)$。当然 $f$ 是一个 $2:1$ 覆叠，在 $f'(x)=0$ 处分歧，即在 $x=-1/2t_2$ 处分歧。所以 $\check X$ 的纤维一般而言是不连通的，而 $W_0=-1/2t_2$ 是 $W$ 的多重纤维。这与我们想要的行为截然不同，被称为野性行为（wild behaviour）。显然当 $t_2\to0$ 时，额外的临界点跑到 $\infty$，这就是为什么在 $t_2=0$ 时我们看不到这种行为。但我们需要避开这类势。
处理这个问题有多种不同途径。我们这里采取的办法是只把 $t_2,\dots,t_n$ 当作形式变量。例如，我们可以取 $\cM=\operatorname{Spec}\C[t_0,t_1][\![t_2,\dots,t_n]\!]$。更符合 A-模型讨论的做法是，我们直接同一个类似于 §2.1.4 中所定义的模空间 $\widetilde{\cM}$ 打交道。

\begin{definition}
定义 $\widetilde{\cM}$ 为环空间 $(\C,\cO_{\widetilde{\cM}})$，其中 $\C$ 看作带坐标 $t_1$ 的复流形，而 $\cO_{\widetilde{\cM}}$ 在开集 $U\subseteq\C$ 上的截面由形式幂级数
\[
\sum f_{i_0i_2\cdots i_n}t_0^{i_0}t_2^{i_2}\cdots t_n^{i_n}
\]
组成，其中 $f_{i_0i_2\cdots i_n}$ 是 $U$ 上的全纯函数。
\end{definition}
% [Tutor Note: p.75 末起 Definition 2.36 的句子跨到 p.76（下页首句“with $f_{i_0i_2\cdots i_n}$ a holomorphic function on $U$”），按 STYLE_GUIDE §10“跨页句子开始者写完”已在本块补全并闭合 definition 环境；c02d 请从 p.76 的“Let $\check X$ be the subspace of $\widetilde{\cM}\times\C^{n+1}$ defined by the equation $e^{t_1}=x_0\cdots x_n$”直接续写，勿重复本句。]
